% Finite-dimensional Gaussian identities used before the path integral.

高斯积分是二次作用量与自由传播核的有限维原型。对普通实变量 $x$ 和 $a>0$，定义
\begin{equation}
 I(a)\equiv\int_{-\infty}^{\infty}\dd x\,\ee^{-ax^2}.
 \label{eq:real-gaussian-definition-dp68}
\end{equation}
这个积分的精确归一化为
\begin{equation}
 \boxed{I(a)=\sqrt{\frac{\pi}{a}}.}
 \label{eq:real-gaussian-r15}
\end{equation}
量子传播与 Fourier 型积分还会遇到纯振荡的二次相位。其积分值由加入正收敛因子后的绝对收敛积分连续延拓定义。在这一规定下，$A>0$ 时
\begin{equation}
 \boxed{
 \int\frac{\dd p}{2\pi\hbar}\,
 \ee^{-\ii A p^2/(2\hbar)}
 =\frac{1}{\sqrt{2\pi\ii\hbar A}}.}
 \label{eq:fresnel-normalized-r15}
\end{equation}
平方根的支由同一个收敛处方固定；传播核中的相位与归一化必须继承这一支选择。
